% =============================================================================
% 1.8 Traitement des documents structurés
% =============================================================================
\section{Traitement des documents structurés}
\label{sec:ch1-documents}

La qualité d'un système RAG est bornée par la qualité de ce qu'il ingère.
Un extrait mal segmenté, un tableau aplati en flux de caractères ou un
article dont le numéro a disparu produisent, en aval, une recherche
imprécise et une génération infidèle. Le traitement des documents
structurés n'est donc pas un préambule technique~: c'est un sujet
d'état de l'art, car il détermine ce qui est ensuite recherchable et
citable.

\subsection{De l'OCR à la compréhension de documents}

La reconnaissance optique de caractères a connu une évolution qui dépasse
largement la transcription de glyphes. Le tableau~\ref{tab:evolution-ocr}
en retrace les générations.

\begin{btptable}{Évolution de la reconnaissance de documents}
\label{tab:evolution-ocr}
\begin{tabularx}{\linewidth}{@{}L{2.6cm} Y Y@{}}
\toprule
\rowcolor{btppale}
\btpth{Période} & \btpth{Approche} & \btpth{Limite principale} \\
\midrule
1960--1990 & Reconnaissance matricielle de caractères imprimés. & Polices variables, manuscrit. \\
1990--2005 & OCR statistique (famille Tesseract). & Perte de la mise en page. \\
2005--2015 & OCR et détection de blocs, colonnes, orientations. & Fragilité sur les documents complexes. \\
2015--2020 & Extraction dédiée de tableaux. & Pas de compréhension sémantique. \\
Depuis 2020 & Compréhension de documents (texte et mise en page conjoints). & Maturité encore en consolidation. \\
\bottomrule
\end{tabularx}
\end{btptable}

La rupture conceptuelle est celle de la \btpterm{compréhension de
documents} (\anglais{Document Understanding}). Là où l'OCR traditionnel
produit un flux linéaire de texte, cette approche vise à préserver la
structure hiérarchique~: titres, sous-titres, articles numérotés, tableaux,
notes de bas de page, annexes. Elle s'appuie sur des modèles de fondation
entraînés conjointement sur le texte et la géométrie de la page
(LayoutLMv3, DocFormer et successeurs). Des plateformes comme Docling,
Azure Document Intelligence ou AWS Textract en sont des illustrations
industrielles.

\begin{btpbox}[Pourquoi la structure est non négociable]
Un CCAP est organisé en articles~; un CCTP suit une hiérarchie de sections~;
un DPGF est un tableau. Perdre cette structure lors de l'extraction
dégrade trois capacités à la fois~: la recherche (un article entier est
plus pertinent qu'un fragment), la citation (la référence à un article
précis est exigée) et la génération (le plan de la réponse s'appuie sur
celui du dossier entrant).
\end{btpbox}

\subsection{Analyse de mise en page et reconnaissance de tableaux}

Les tableaux constituent une composante critique des dossiers d'appels
d'offres. Un bordereau de prix ou une décomposition du prix global
contiennent des informations financières dont l'altération compromet
toute génération ultérieure de sections chiffrées.

Trois défis se distinguent~:

\begin{itemize}
  \item \btpterm{Reconnaissance des cellules}~: identifier les en-têtes, les
        lignes, les colonnes et les cellules fusionnées~;
  \item \btpterm{Préservation des hiérarchies}~: certains tableaux
        embarquent des sous-totaux ou des regroupements par lot~;
  \item \btpterm{OCR spécifique}~: les tableaux emploient souvent des
        polices plus petites et moins régulières que le corps du texte.
\end{itemize}

\begin{btptable}{Approches de reconnaissance de tableaux}
\label{tab:tableaux}
\begin{tabularx}{\linewidth}{@{}L{3.4cm} Y Y@{}}
\toprule
\rowcolor{btppale}
\btpth{Approche} & \btpth{Avantages} & \btpth{Limites} \\
\midrule
Règles heuristiques (lignes, colonnes) & Rapide, simple. & Fragile dès que le quadrillage est irrégulier. \\
Détection par réseau convolutif & Plus robuste aux variations de mise en page. & Exige des données d'entraînement. \\
Modèles de fondation multimodaux & Très robuste, comprend le rôle des cellules. & Coût computationnel élevé. \\
\bottomrule
\end{tabularx}
\end{btptable}

\subsection{Enjeux pour les appels d'offres}

Les dossiers d'appels d'offres camerounais présentent des caractéristiques
qui aggravent ces défis. Une part non négligeable des pièces circule encore
sous forme de scans, parfois de photocopies de mauvaise qualité, portant
tampons, signatures et annotations manuscrites. Un pipeline unique, conçu
pour des PDF nativement numériques, échoue sur ce corpus.

Il en résulte une exigence de \btpterm{détection en amont}~: distinguer un
document numérique, pour lequel une analyse structurelle suffit, d'un
document numérisé, pour lequel une reconnaissance optique avec analyse de
mise en page s'impose. Le choix des outils concrets relève du
chapitre~\ref{chap:resultats}~; le constat d'état de l'art est que cette
bifurcation n'est plus optionnelle.
